% Generated draft: requires mathematical and source review.
\documentclass[11pt]{article}
\usepackage[margin=25mm]{geometry}
\usepackage{fontspec}
\setmainfont{Latin Modern Roman}
\usepackage{amsmath,amssymb,mathtools,unicode-math}
\setmathfont{Latin Modern Math}
\usepackage{enumitem,needspace,xurl,hyperref}
\setlength{\parindent}{0pt}
\setlength{\parskip}{5pt}
\setlength{\emergencystretch}{3em}
\newcommand{\N}{\mathbb{N}}
\newcommand{\Z}{\mathbb{Z}}
\newcommand{\Q}{\mathbb{Q}}
\newcommand{\R}{\mathbb{R}}
\newcommand{\C}{\mathbb{C}}
\newcommand{\F}{\mathbb{F}}
\newcommand{\A}{\mathbb{A}}
\newcommand{\PP}{\mathbb{P}}
\newcommand{\E}{\mathbb{E}}
\newcommand{\eps}{\varepsilon}
\newcommand{\dd}{\,\mathrm{d}}
\begin{document}
\section*{UnsolvedMath ID: 700740}
\textbf{Source problem code:} catalogue entry
\subsection{Definitions and mathematical statement}
Let \(F=F(x,y)\) be the free group on two generators \(x,y\). For this draft, use the commutator convention \([a,b]=a^{-1}b^{-1}ab\) and interpret iterated commutators as left-normed. Thus, define \(w_0=x\) and \(w_{j+1}=[w_j,y]\) for \(0\leq j<6\), so that
\[
w_6=[x,y,y,y,y,y,y].
\]
The supplied record does not explicitly specify these commutator conventions; they require confirmation against the cited source.

Let
\[
N_5=\left\langle g^5\mid g\in F\right\rangle,
\qquad B(2,5)=F/N_5.
\]
Here \(N_5\) is a normal subgroup, and \(B(2,5)\) is the free Burnside group of rank two and exponent dividing five. Since \((g^5)^{-1}=(g^{-1})^5\), membership in \(N_5\) means being a finite product of fifth powers.

The question is whether
\[
w_6\in N_5,
\]
or, equivalently, whether there exist an integer \(m\geq0\) and elements \(g_1,\ldots,g_m\in F\) such that
\[
[x,y,y,y,y,y,y]=g_1^5\cdots g_m^5,
\]
where the empty product is the identity. The source additionally states that a negative answer would imply that \(B(2,5)\) is infinite. This conditional consequence is preserved as a supplied source assertion.
\subsection{Short English statement}
Is the left-normed commutator of \(x\) with six successive copies of \(y\) a finite product of fifth powers in the free group on \(x,y\)? The source states that a negative answer would imply that \(B(2,5)\) is infinite. The commutator conventions used here require source confirmation.
\subsection{Sources}
\begin{itemize}
\item Supplied local UnsolvedMath registry record: permanent ID \(700740\), exact source problem code \texttt{catalogue entry}, title \emph{Kourovka Notebook Problem 11.48}; upstream revision \texttt{d371ee917ac9647fbb06cd6fc529deec6f3dda82}.
\item Attributed to A. I. Kostrikin, proposed in 1990. Supplied bibliographic source: \emph{The Kourovka Notebook (Unsolved Problems in Group Theory)}, 11th issue (1990), Problem 11.48; also identified with the 21st issue edition (2026), \texttt{arXiv:1401.0300v48}. The supplied collection description names E. I. Khukhro and V. D. Mazurov as editors.
\item The local record labels the problem \texttt{open} and records an update date of 2026-10-16. These are supplied registry metadata, not an independently verified mathematical status or literature assessment. No independent source consultation is claimed.
\end{itemize}
\end{document}
